\section{Specialized Architectures}\label{sec:deep_learning_specialized_architectures}

The concepts presented so far have been introduced with the \gls{mlp}.
In this architecture, each neuron is connected to all neurons of the previous layer, and no assumption is made about the structure of the input.
However, many types of data have a structure that can be exploited.
The pixels of an image are organized in a spatial grid, and the words of a sentence form an ordered sequence in which each word depends on its context.
Specialized architectures organize the connections between neurons so as to take this structure into account.

These architectures do not modify the principles introduced in this chapter.
A network remains a composition of layers, where each layer computes its representation from that of the previous one,
\begin{equation}
\latentRepresentation[l] = \LayerFunction{l}\left(\latentRepresentation[l-1]\right).
\label{eq:layer_function}
\end{equation}
Its parameters are still optimized by gradient descent, and it is still evaluated on data it has not seen during training.
Only the function $\LayerFunction{l}$ computed by each layer changes.
In particular, each layer still produces an intermediate representation of the input, which may take the form of feature maps or of a sequence of tokens.
This representation is shaped by the same learning process as in a \gls{mlp}.
This section presents two architectures used throughout this thesis, the \gls{cnn} and the attention mechanism, which is the basis of the Transformer architecture.

\subsection{Convolutional Neural Networks}\label{subsec:deep_learning_convolutional_neural_networks}

\glspl{cnn} were introduced to process images \parencite{lecun1989backpropagation}.
Applying a \gls{mlp} to an image would require connecting each neuron to every pixel.
The number of parameters would then grow with the size of the image, and the spatial organization of the pixels would be ignored.
\glspl{cnn} rely instead on two properties of images.
The first is locality, since the information needed to detect a pattern such as an edge is contained in a small region of the image.
The second is that the same pattern can appear at any position in the image.

A convolutional layer exploits these two properties by means of kernels, also called filters.
A kernel is a small matrix of weights, typically of size $3 \times 3$, which is applied to a local region of the representation of the previous layer, called its receptive field.
The kernel is then slid over every position of this representation.
At each position, a weighted sum of the elements of the receptive field is computed, which gives one element of the output.
The set of outputs obtained over all positions forms a feature map.
Since the same weights are used at every position, the number of parameters of the layer does not depend on the size of the input, and a pattern is detected in the same way wherever it appears.

A convolutional layer uses several kernels, each detecting a different pattern.
The representation of layer $l$ is therefore composed of $n_l$ feature maps, also called channels, and is denoted $\mat{H}^{(l)}$.
The feature map produced by the $k$-th kernel is given by:
\begin{equation}
\mat{H}^{(l)}_k = \sigma^{(l)}\left(\mat{W}^{(l)}_k \ast \mat{H}^{(l-1)} + b^{(l)}_k\right), \qquad k = 1, \dots, n_l,
\label{eq:convolutional_layer}
\end{equation}
where:
\begin{itemize}
\item $\mat{W}^{(l)}_k$ is the $k$-th kernel of layer $l$,
\item $b^{(l)}_k$ is the bias associated with this kernel,
\item $\ast$ denotes the convolution operation, which applies the kernel at every position of $\mat{H}^{(l-1)}$.
\end{itemize}
For a single-channel input and a kernel of size $\kappa \times \kappa$, the element at position $(i, j)$ of the convolution is:
\begin{equation}
\left(\mat{W}^{(l)}_k \ast \mat{H}^{(l-1)}\right)_{i,j} = \sum_{i'=1}^{\kappa} \sum_{j'=1}^{\kappa} W^{(l)}_{k,i',j'} \, H^{(l-1)}_{i+i'-1,\, j+j'-1}.
\label{eq:convolution}
\end{equation}
When $\mat{H}^{(l-1)}$ has several channels, each kernel covers all of them, and the sum also runs over the $n_{l-1}$ channels.
For the first layer, $\mat{H}^{(0)}$ is the input image, whose channels correspond for instance to the color components of the pixels.
The kernel can also be moved by steps larger than one element, called the stride, which reduces the size of the feature maps.

Figure~\ref{fig:convolutional_layer} illustrates this computation for a single-channel input.
The receptive field, highlighted in orange on the representation $\mat{H}^{(l-1)}$, is combined with each kernel $\mat{W}^{(l)}_k$.
Each combination gives one element of the corresponding feature map $\mat{H}^{(l)}_k$, shown in the same color as the kernel.
Sliding the receptive field over every position of $\mat{H}^{(l-1)}$ produces all the other elements of these feature maps.

\begin{figure}[t]
    \centering
    \resizebox{0.85\textwidth}{!}{
        \begin{tikzpicture}

        %%── Input feature map (layer l-1): 6x6 grid, cell = 0.5 ─────────────────────
        \begin{scope}[shift={(0,-1.5)}]
            \fill[gray!10] (0,0) rectangle (3,3);
            \fill[myorange!35] (0.5,1.0) rectangle (2.0,2.5);
            \foreach \i in {0,...,6} {
                \draw[gray!60] (0.5*\i,0) -- (0.5*\i,3);
                \draw[gray!60] (0,0.5*\i) -- (3,0.5*\i);
            }
            \draw[thick] (0,0) rectangle (3,3);
            \draw[very thick, myorange] (0.5,1.0) rectangle (2.0,2.5);
        \end{scope}
        \node at (1.5,2.0) {Layer $l-1$};
        \node[anchor=east] at (-0.1,0) {$\mat{H}^{(l-1)}$};
        \node[text=myorange] at (1.5,-2.0) {Receptive field};

        %%── Column titles ───────────────────────────────────────────────────────────
        \node at (5.8,4.9) {Kernels};
        \node at (10,4.9) {Layer $l$};

        %%── One kernel and its output feature map per row ───────────────────────────
        \foreach \yc/\col/\k in {3.2/myblue/1, 0/mygreen/2, -3.2/myred/n_l} {

            % Kernel: 3x3 grid, cell = 0.4
            \begin{scope}[shift={(5.2,\yc-0.6)}]
                \fill[\col!25] (0,0) rectangle (1.2,1.2);
                \foreach \i in {0,...,3} {
                    \draw[\col] (0.4*\i,0) -- (0.4*\i,1.2);
                    \draw[\col] (0,0.4*\i) -- (1.2,0.4*\i);
                }
                \draw[thick, \col] (0,0) rectangle (1.2,1.2);
                \node[text=\col] at (0.6,1.5) {$\mat{W}^{(l)}_{\k}$};
            \end{scope}

            % Output feature map: 4x4 grid, cell = 0.5
            \begin{scope}[shift={(9,\yc-1)}]
                \fill[\col!12] (0,0) rectangle (2,2);
                \fill[\col!70] (0.5,1.0) rectangle (1.0,1.5);
                \foreach \i in {0,...,4} {
                    \draw[\col!50] (0.5*\i,0) -- (0.5*\i,2);
                    \draw[\col!50] (0,0.5*\i) -- (2,0.5*\i);
                }
                \draw[thick, \col] (0,0) rectangle (2,2);
                \node[text=\col, anchor=west] at (2.15,1) {$\mat{H}^{(l)}_{\k}$};
            \end{scope}

            % Receptive field -> kernel
            \draw[dashed, \col] (2.0, 1.0) -- (5.2,\yc+0.6);
            \draw[dashed, \col] (2.0,-0.5) -- (5.2,\yc-0.6);
            % Kernel -> output cell
            \draw[dashed, \col] (6.4,\yc+0.6) -- (9.5,\yc+0.5);
            \draw[dashed, \col] (6.4,\yc-0.6) -- (9.5,\yc);
        }

        %%── Vertical dots between the second and the last kernel ───────────────────
        \node at (5.8,-1.6) {$\vdots$};
        \node at (10,-1.6) {$\vdots$};

        \end{tikzpicture}
    }
    \caption{Convolutional layer.}
    \label{fig:convolutional_layer}
\end{figure}


As with the \gls{mlp}, convolutional layers are stacked.
The receptive field of a neuron, measured on the input image, grows with depth.
The first layers thus detect simple local patterns, while the deeper layers combine them into patterns that cover larger regions of the image \parencite{goodfellow2016deep}.
The last feature maps are then flattened into a vector and passed through a \gls{mlp}, which produces the output of the network.
More generally, the feature maps $\mat{H}^{(l)}$ of any convolutional layer can be flattened into a vector, which plays the same role as the representation $\latentRepresentation[l]$ of a layer of a \gls{mlp}.
\glspl{cnn} trained on \glspl{gpu} were at the origin of the rise of \gls{dl} in computer vision \parencite{krizhevsky2012imagenet}.
This combination of convolutional layers followed by a \gls{mlp} is also widely used in \gls{rl} to process visual observations \parencite{mnih2013playingatarideepreinforcement, espeholt2018impala}.

\subsection{Attention Mechanism}\label{subsec:deep_learning_attention_mechanism}

The attention mechanism processes an input made of a sequence of elements, called tokens.
Depending on the task, a token can represent a part of a word, a region of an image, or an element of an observation.
Each token is a vector, and the $m$ tokens of layer $l-1$ are stacked as the rows of a matrix $\mat{H}^{(l-1)}$.
Unlike a convolutional layer, which only combines neighboring elements, the attention mechanism allows each token to gather information from all the other tokens of the sequence, whatever their position \parencite{vaswani2017attention}.

Three matrices are first computed from the tokens by learned linear projections:
\begin{equation}
\mat{Q} = \mat{H}^{(l-1)} \mat{W}^{(l)}_Q, \qquad
\mat{K} = \mat{H}^{(l-1)} \mat{W}^{(l)}_K, \qquad
\mat{V} = \mat{H}^{(l-1)} \mat{W}^{(l)}_V,
\label{eq:queries_keys_values}
\end{equation}
where the rows of $\mat{Q}$, $\mat{K}$ and $\mat{V}$ are respectively the queries, the keys and the values of the tokens.
The query of a token describes the information it is looking for, its key describes the information it contains, and its value is the information it transmits.
The queries and keys are vectors of dimension $\KeyDim$.

The query of each token is then compared with the keys of all tokens by a scalar product.
These scores are normalized into attention weights by the softmax function, applied to each row:
\begin{equation}
\mat{A} = \softmax\left(\frac{\mat{Q}\mat{K}^{\top}}{\sqrt{\KeyDim}}\right),
\qquad \text{with} \quad \softmax(\vect{u})_i = \frac{\exp(u_i)}{\sum_{j} \exp(u_j)}.
\label{eq:attention_weights}
\end{equation}
The weight $A_{i,j}$ indicates how much token $i$ attends to token $j$.
The weights of each row are positive and sum to one.
The division by $\sqrt{\KeyDim}$ prevents the scores from becoming too large as the dimension increases, which would concentrate the softmax on a single token.
Finally, each output token is the weighted sum of the values according to these weights:
\begin{equation}
\mat{H}^{(l)} = \mat{A}\mat{V} = \softmax\left(\frac{\mat{Q}\mat{K}^{\top}}{\sqrt{\KeyDim}}\right) \mat{V}.
\label{eq:attention}
\end{equation}

Figure~\ref{fig:attention_mechanism} illustrates this mechanism for a sequence of $m = 4$ tokens, numbered from $1$ to $4$ and stacked as the rows of $\mat{H}^{(l-1)}$.
The tokens are projected into queries, keys and values, shown in blue, green and red respectively.
The queries and keys are combined to give the attention weights $\mat{A}$, represented in orange as a heat map in which darker cells correspond to larger weights.
Row $i$ of this heat map indicates how much token $i$ attends to each token $j$.
The output tokens of layer $l$, shown in purple, are then obtained by weighting the values with these attention weights.

\begin{figure}[t]
    \centering

    % \matgrid{x}{y}{cols}{rows}{color}: grid with lower-left corner at (x,y), cell = 0.4.
    \newcommand{\matgrid}[5]{%
        \begin{scope}[shift={(#1,#2)}]
            \fill[#5!20] (0,0) rectangle (0.4*#3,0.4*#4);
            \foreach \i in {0,...,#3} \draw[#5!70] (0.4*\i,0) -- (0.4*\i,0.4*#4);
            \foreach \j in {0,...,#4} \draw[#5!70] (0,0.4*\j) -- (0.4*#3,0.4*\j);
            \draw[thick, #5] (0,0) rectangle (0.4*#3,0.4*#4);
        \end{scope}}

    \resizebox{\textwidth}{!}{
        \begin{tikzpicture}[
            arr/.style={-{Stealth[length=4pt, width=3pt]}, thick},
            op/.style={draw, thick, rounded corners=3pt, fill=gray!10, inner sep=5pt},
        ]

        %%── Input tokens H^{(l-1)} ──────────────────────────────────────────────────
        \matgrid{0}{-0.8}{4}{4}{darkgray}
        \foreach \r in {1,...,4}
            \node[anchor=east, font=\footnotesize] at (-0.05,1.0-0.4*\r) {$\r$};
        \node at (0.8,1.3) {Layer $l-1$};
        \node at (0.8,-1.3) {$\mat{H}^{(l-1)}$};

        %%── Queries, keys and values ────────────────────────────────────────────────
        \matgrid{3.6}{1.8}{3}{4}{myblue}
        \matgrid{3.6}{-0.8}{3}{4}{mygreen}
        \matgrid{3.6}{-3.4}{3}{4}{myred}
        \node[text=myblue]  at (4.2,3.7)  {Queries $\mat{Q}$};
        \node[text=mygreen] at (4.2,1.1)  {Keys $\mat{K}$};
        \node[text=myred]   at (4.2,-1.5) {Values $\mat{V}$};

        \draw[arr, myblue]  (1.7,0.3)  -- node[fill=white, inner sep=1.5pt] {$\mat{W}^{(l)}_Q$} (3.5,2.6);
        \draw[arr, mygreen] (1.7,0)    -- node[fill=white, inner sep=1.5pt] {$\mat{W}^{(l)}_K$} (3.5,0);
        \draw[arr, myred]   (1.7,-0.3) -- node[fill=white, inner sep=1.5pt] {$\mat{W}^{(l)}_V$} (3.5,-2.6);

        %%── Scores and softmax ──────────────────────────────────────────────────────
        \node[op] (score) at (6.7,1.3) {$\dfrac{\mat{Q}\mat{K}^{\top}}{\sqrt{\KeyDim}}$};
        \node[op] (soft)  at (8.7,1.3) {$\softmax$};
        \draw[arr, myblue]  (4.9,2.6) -- (score);
        \draw[arr, mygreen] (4.9,0)   -- (score);
        \draw[arr] (score) -- (soft);

        %%── Attention weights A (heat map) ──────────────────────────────────────────
        \begin{scope}[shift={(10.9,0.5)}]
            \foreach \r/\a/\b/\c/\d in {0/70/15/10/5, 1/20/50/25/5, 2/5/35/40/20, 3/5/10/30/55} {
                \fill[myorange!\a] (0.0,1.2-0.4*\r) rectangle (0.4,1.6-0.4*\r);
                \fill[myorange!\b] (0.4,1.2-0.4*\r) rectangle (0.8,1.6-0.4*\r);
                \fill[myorange!\c] (0.8,1.2-0.4*\r) rectangle (1.2,1.6-0.4*\r);
                \fill[myorange!\d] (1.2,1.2-0.4*\r) rectangle (1.6,1.6-0.4*\r);
            }
            \foreach \i in {0,...,4} {
                \draw[myorange!50] (0.4*\i,0) -- (0.4*\i,1.6);
                \draw[myorange!50] (0,0.4*\i) -- (1.6,0.4*\i);
            }
            \draw[thick, myorange] (0,0) rectangle (1.6,1.6);
            \foreach \r in {1,...,4} {
                \node[anchor=east, font=\footnotesize] at (-0.05,1.8-0.4*\r) {$\r$};
                \node[anchor=south, font=\footnotesize] at (-0.2+0.4*\r,1.62) {$\r$};
            }
            \node[font=\footnotesize] at (-0.55,0.8) {$i$};
            \node[font=\footnotesize] at (0.8,2.15) {$j$};
        \end{scope}
        \node[text=myorange] at (11.7,3.1) {Attention weights $\mat{A}$};
        \draw[arr] (soft) -- (9.9,1.3);

        %%── Weighted sum of the values ──────────────────────────────────────────────
        \node[op] (mul) at (11.7,-2.6) {$\mat{A}\mat{V}$};
        \draw[arr, myorange] (11.7,0.4) -- (mul);
        \draw[arr, myred] (4.9,-2.6) -- (mul);

        %%── Output H^{(l)} ──────────────────────────────────────────────────────────
        \matgrid{13.6}{-3.4}{3}{4}{mypurple}
        \draw[arr] (mul) -- (13.5,-2.6);
        \foreach \r in {1,...,4}
            \node[anchor=west, font=\footnotesize, text=mypurple] at (14.85,-1.6-0.4*\r) {$\r$};
        \node[text=mypurple] at (14.2,-0.9) {Layer $l$};
        \node[text=mypurple] at (14.2,-3.9) {$\mat{H}^{(l)}$};

        \end{tikzpicture}
    }
    \caption{Attention mechanism.}
    \label{fig:attention_mechanism}
\end{figure}


The main difference with the layers presented so far lies in the weights that combine the inputs.
In a \gls{mlp} or a \gls{cnn}, these weights are parameters, fixed once training is completed.
In the attention mechanism, the weights $\mat{A}$ are computed from the input itself, so that the way tokens are combined adapts to each input.
Only the projection matrices $\mat{W}^{(l)}_Q$, $\mat{W}^{(l)}_K$ and $\mat{W}^{(l)}_V$ are learned, and their size does not depend on the number of tokens $m$.
Since this mechanism does not take the order of the tokens into account, information about the position of each token is generally added to its vector before the first layer.

The Transformer architecture is built by stacking blocks that combine an attention mechanism with a \gls{mlp} applied independently to each token \parencite{vaswani2017attention}.
In practice, several attention mechanisms are computed in parallel within each block, each with its own projections, which allows different relations between tokens to be captured simultaneously.
This architecture underlies \glspl{llm}, which are trained on very large text corpora to predict the next token of a sequence \parencite{brown2020language}.
In \gls{rl}, it is used to build policies that retain a memory of past observations \parencite{parisotto2020stabilizing}.
Because the attention weights $\mat{A}$ indicate which tokens contribute to each output, they are also used to explain the decisions of agents, as discussed in Section~\ref{subsubsec:attention_based_architecture}.

Whatever the architecture, a network therefore produces a sequence of intermediate representations, whether they take the form of activation vectors, feature maps, or sequences of tokens.
These representations are shaped by the same learning process, since the parameters of every layer are adjusted by gradient descent to minimize the same loss.
Methods that analyze the representations of a network layer by layer can therefore be applied regardless of its architecture.
In particular, the interpretability studies of \glspl{llm} presented in the following chapter rely on the representations produced by the layers of Transformers.
